\chapter{第十一届大唐杯研究生组省赛}
\fancyhead[L]{\color{H6}\kaishu\faIcon{atom}\;第十一届大唐杯研究生组4月19日省赛}
\fancyhead[R]{\color{H6}\kaishu\rightmark\,}

\date{2024年4月19日}{}{\href{https://github.com/xubenshan/dtcup}{\textbf{才疏学浅的小熊}}}
{}
{https://www.xiaohongshu.com/user/profile/6192219d0000000021028647}{小红书}
{https://space.bilibili.com/692546609}{B站账号}



\section{单选题（共40小题，共120分）}



\begin{choice}{A}[]
    5G系统核心网元SMF的功能不包括
    \begin{tasks}(2)
        \task NAS信令及信令的加密和完整性保护
        \task UE IP地址的分配与管理
        \task 合法监听
        \task 会话的建立修改删除
    \end{tasks}
\end{choice}


\begin{choice}{A}[]
    AI大数据模型在移动通信网络中的多重作用中不包括
    \begin{tasks}(4)
        \task 网络部署
        \task 安全保障
        \task 无线资源管理
        \task 用户行为分析
    \end{tasks}
\end{choice}




\begin{choice}{B}[]
    以下选项不属于Gradient Boosting学习算法的是
    \begin{tasks}(4)
        \task Xgboost
        \task cart树
        \task GBDT
        \task lightGBM
    \end{tasks}
\end{choice}



\begin{choice}{C}[]
    衡量系统无线网络覆盖率的重要指标是
    \begin{tasks}(4)
        \task SINR
        \task RSRQ
        \task RSRP
        \task RSSI
    \end{tasks}
\end{choice}


\begin{choice}{B}[]
    电信网的传输设备根据传输介质的不同分为：光纤传输设备、卫星传输设备、（\quad）、缆线传输设备等。
    \begin{tasks}(4)
        \task 业务传输设备
        \task 无线传输设备
        \task 双绞线传输设备
        \task 交换传输设备
    \end{tasks}
\end{choice}

\begin{choice}{D}[]
    以下不属于6G网络高精度通信目标的是
    \begin{tasks}(4)
        \task 协同保证
        \task 吞吐量保证
        \task 无损网络
        \task 全息网络
    \end{tasks}
\end{choice}




\begin{choice}{A}[]
    IPMT流程第一阶段的输出成果是什么
    \begin{tasks}(2)
        \task 新产品规划建设报告
        \task 新产品
        \task 客户反馈信息
        \task 企业自身能力
    \end{tasks}
\end{choice}


\begin{choice}{D}[]
    以下计算公式错误的是
    \begin{tasks}(1)
        \task $XdB-YdB=(X-Y)dB$
        \task $XdBm+YdB=(X+Y)dBm$
        \task $XdB+YdB=(X+Y)dB$
        \task $XdBm+YdBm=(X+Y)dBm$
    \end{tasks}
\end{choice}


\begin{choice}{C}[]
    卫星地球站通过赤道上空约36000km的通信卫星的转发进行通信，视地球站纬度高低，其一跳的单程空间距离为72000-（\quad）km。
    \begin{tasks}(4)
        \task 90000
        \task 92000
        \task 80000
        \task 82000
    \end{tasks}
\end{choice}


\begin{choice}{B}[]
    大规模天线中，上行支持的MIMO方式有
    \begin{tasks}(2)
        \task 基于CRS的MU-MIMO
        \task 基于DMRS的SU-MIMO
        \task 不支持波束赋型发送
        \task 基于CSI-RS的MU-MIMO
    \end{tasks}
\end{choice}

\begin{choice}{A}[]
    下面哪个信令可以携带SCG Add或PSCell Change相关SCG配置
    \begin{tasks}(1)
        \task LTE侧的RRC Connection Reconfiguration
        \task NR侧的RRC Connection Setup
        \task LTE侧的RRC Connection Setup
        \task NR侧的RRC Reconfiguration
    \end{tasks}
\end{choice}


\begin{choice}{A}[]
    以下哪个选项不是5G接入网特征
    \begin{tasks}(2)
        \task 多个终端地址
        \task 多种接入技术
        \task 多层异构
        \task 多场景
    \end{tasks}
\end{choice}


\begin{choice}{D}[]
    IPMT流程第二阶段的输出成果是
    \begin{tasks}(2)
        \task 新产品
        \task 战略投资决策
        \task 产品测试报告
        \task 新产品立项批准报告
    \end{tasks}
\end{choice}


\begin{choice}{A}[]
    6G网络将采用一种由大量超材料单元构成的无源阵列，能以很低的功耗和成本灵活调控无线环境中的电磁波，从而大幅提高接受信号质量的关键技术，该技术是
    \begin{tasks}(2)
        \task 智能超表面技术
        \task 电磁波角动能
        \task 无线蜂窝大规模MIMO技术
        \task 无源天线阵列技术
    \end{tasks}
\end{choice}

\begin{choice}{A}[]
    在5G网络规划PCI时，应考虑PCI的mod（\quad）影响
    \begin{tasks}(4)
        \task 3
        \task 6
        \task 4
        \task 2
    \end{tasks}
\end{choice}

\begin{choice}{A}[]
    以下说法错误的是
    \begin{tasks}(1)
        \task 产品的变动成本仅存在于生产的作业环节中
        \task 产品的全成本不是在一个作业环节中形成的
        \task 成本控制以完成预定成本限额为目标
        \task 产品策划阶段最重要的工作就是对新产品的正确描述
    \end{tasks}
\end{choice}



\begin{choice}{A}[]
    激光雷达以激光作为载波，激光是光波段电磁辐射，波长比微波和毫米波
    \begin{tasks}(4)
        \task 短
        \task 以上均不对
        \task 一样长
        \task 长
    \end{tasks}
\end{choice}

\begin{choice}{B}[]
    5G NR中，开启异频测量的事件是
    \begin{tasks}(4)
        \task A3
        \task A2
        \task A1
        \task B1
    \end{tasks}
\end{choice}


\begin{choice}{A}[]
    5G中：FR2频段通信的优势不包括
    \begin{tasks}(2)
        \task 穿透能力强
        \task 波束集中，提高能量
        \task 可用频段宽
        \task 方向性好，受干扰影响小
    \end{tasks}
\end{choice}



\begin{choice}{A}[]
    IPMT职责不包括哪个
    \begin{tasks}(2)
        \task 产品销售
        \task 战略投资决策
        \task 产品布局
        \task 竞争策略制定
    \end{tasks}
\end{choice}



\begin{choice}{D}[]
    CAN总线使用的数据编码是
    \begin{tasks}(2)
        \task 归零码
        \task 差分曼彻斯特编码
        \task 曼彻斯特编码
        \task 非归零码
    \end{tasks}
\end{choice}




\begin{choice}{A}[]
    在5G NR中，一个REG包含(\quad)个RB
    \begin{tasks}(4)
        \task 1
        \task 2
        \task 3
        \task 4
    \end{tasks}
\end{choice}

\begin{choice}{C}[]
    5G NR有多少个唯一的物理层小区ID
    \begin{tasks}(4)
        \task 504
        \task 512
        \task 1008
        \task 256
    \end{tasks}
\end{choice}

\begin{choice}{B}[]
    在实际应用中，GPS接收装置利用几颗以上卫星信号来定出使用者所在位置
    \begin{tasks}(4)
        \task 2
        \task 4
        \task 6
        \task 1
    \end{tasks}
\end{choice}


\begin{choice}{A}[]
    华为Mate 60pro的卫星通话功能采用中国自主研制的天通一号卫星移动通信系统，天通一号卫星是（\quad）同步卫星
    \begin{tasks}(2)
        \task 高轨道
        \task 低轨道
        \task 中高轨道
        \task 近地轨道
    \end{tasks}
\end{choice}


\begin{choice}{D}[]
    下列不属于5GmMTC应用场景的是
    \begin{tasks}(4)
        \task 智能抄表
        \task 智慧城市
        \task 智慧家居
        \task 高清视频
    \end{tasks}
\end{choice}



\begin{choice}{B}[]
    对于5G支持的频段，说法错误的是
    \begin{tasks}(1)
        \task FR1频段下的最大带宽为273RB
        \task FR1频段下，任何频带都能支持100MHz带宽
        \task FR2频段下的最大带宽为264RB
        \task FR2频段下只有120khz的子载波带宽能够支持400MHz小区带宽
    \end{tasks}
\end{choice}


\begin{choice}{C}[]
    下面哪个原因不是采用异构网络的理由
    \begin{tasks}(2)
        \task 深度覆盖需要
        \task 热点流量需要
        \task 站址获取难
        \task 宏站覆盖的范围大，上行功率受限
    \end{tasks}
\end{choice}


\begin{choice}{B}[]
    6G网络架构中的四层不包括
    \begin{tasks}(2)
        \task 开放使能层
        \task 传输与接入层
        \task 资源与算力层
        \task 服务化功能层
    \end{tasks}
\end{choice}


\begin{choice}{A}[]
    （\quad）拓扑结构存在若干节点中心，且任何两个节点之间的通信都要经过中心节点
    \begin{tasks}(4)
        \task 树型
        \task 网状型
        \task 环形
        \task 星型
    \end{tasks}
\end{choice}


\begin{choice}{D}[]
    无线网络规划流程的顺序是
    \begin{tasks}(1)
        \task 数据采集-站址规划-无线网络规模估算-参数详细规划-无线网络仿真
        \task 无线网络规模估算-站址规划-无线网络仿真-参数详细规划-数据采集
        \task 站址规划-无线网络仿真-参数详细规划-数据采集-无线网络规模估算
        \task 数据采集-无线网络规模估算-站址规划-无线网络仿真-参数详细规划
    \end{tasks}
\end{choice}


\begin{choice}{D}[]
    bagging和pasting的区别是
    \begin{tasks}(2)
        \task 预测器不同
        \task bagging采样不放回，pasting采样放回
        \task 二者无区别
        \task bagging采样放回，pasting采样不放回
    \end{tasks}
\end{choice}



\begin{choice}{C}[]
    在5G NR中，PDCCH采用的调制方式是
    \begin{tasks}(2)
        \task 16QAM
        \task 256QAM
        \task QPSK
        \task 64QAM
    \end{tasks}
\end{choice}


\begin{choice}{B}[]
    CART算法在树生成过程中，使用基尼指数选择最优特征，生成的树是
    \begin{tasks}(2)
        \task b树
        \task 二叉树
        \task 以上都不是
        \task 红黑树
    \end{tasks}
\end{choice}


\begin{choice}{D}[]
    卫星移动通信业务使用以下哪个频段
    \begin{tasks}(4)
        \task Ku
        \task S
        \task C
        \task L
    \end{tasks}
\end{choice}


\begin{choice}{A}[]
    NR系统中，随机接入的目的错误的是
    \begin{tasks}(2)
        \task 实现UE和gNB之间的下行同步
        \task 获取MSG3的资源
        \task 获取C-RNTI基站识别UE的标识
        \task 实现UE和gNB之间的上行同步
    \end{tasks}
\end{choice}


\begin{choice}{B}[]
    R15版本，5G NR的无线帧长度为
    \begin{tasks}(4)
        \task 20ms
        \task 10ms
        \task 0.5ms
        \task 1ms
    \end{tasks}
\end{choice}


\begin{choice}{B}[]
    对于1个服务小区，基站可以通过专用RRC信令给UE配置多个DLBWP和多个ULBWP，最多各配多少个
    \begin{tasks}(4)
        \task 8
        \task 4
        \task 12
        \task 2
    \end{tasks}
\end{choice}


\begin{choice}{B}[]
    下列哪个因素不影响AAU正常接入
    \begin{tasks}(2)
        \task 对应的HBPOD板卡状态为未初始化状态
        \task 修改了小区物理ID列表
        \task DCPD给AAU供电的开关未闭合
        \task 对应的光模块插反了
    \end{tasks}
\end{choice}


\begin{choice}{C}[]
    降低生产作业环节变动成本的方式不包括
    \begin{tasks}(2)
        \task 整机自动测试
        \task 整机测试生产线柔性化改造
        \task 单板测试生产线刚性化改造
        \task 单板自动测试
    \end{tasks}
\end{choice}



\section{多选题（共30小题，共120分）}

\begin{choice}{\;BD\;}[]
    6G上游行业包括
    \begin{tasks}(2)
        \task 智能制造
        \task 光纤光缆
        \task 自动驾驶
        \task 芯片
    \end{tasks}
\end{choice}

\begin{choice}{\;ABD\;}[]
    5G NR中，可用作下行参考信号包括
    \begin{tasks}(4)
        \task CSI-RS
        \task DMRS
        \task SRS
        \task PT-RS
    \end{tasks}
\end{choice}

\begin{choice}{\;ABD\;}[]
    以下属于高精度地图的作用的有
    \begin{tasks}(1)
        \task 帮助车辆识别车道的结确中心线
        \task 导航地图智能达到米级精度，而高精度地图可以达到厘米级精度
        \task 帮助车辆识别行人
        \task 可以帮助传感器缩小检测范围
    \end{tasks}
\end{choice}

\begin{choice}{\;ABCD\;}[]
    AI大数据模型在移动通信网络中的多重作用包括
    \begin{tasks}(2)
        \task 安全保障
        \task 无线资源管理
        \task 网络优化
        \task 用户行为分析
    \end{tasks}
\end{choice}


\begin{choice}{\;ABCD\;}[]
    5G R15版本，5G核心网包含的功能服务模块包括
    \begin{tasks}(4)
        \task SMF
        \task AMF
        \task UPF
        \task NSSF
    \end{tasks}
\end{choice}


\begin{choice}{\;ABD\;}[]
    NR系统中，PUCCH中UCI携带的信息有
    \begin{tasks}(2)
        \task ACK/NACK
        \task 调度请求(SchedulingRequest)
        \task pointA
        \task CSI(ChannelStateInformation)
    \end{tasks}
\end{choice}

\begin{choice}{\;ACD\;}[]
    变动成本法的优点包括
    \begin{tasks}(1)
        \task 便于分清各部门的经济责任，有利于对成本的控制与业绩的评价
        \task 作为财务会计成本核算方式，已经纳入会计准则，成为通行标准
        \task 能直接提供各种产品的变动成本，可依此直接进行分产品型号的盈亏核算，有利于管理人员的决策分析
        \task 简化了产品成本计算
    \end{tasks}
\end{choice}


\begin{choice}{\;AD\;}[]
    以下选项中属于Boosting提升法的选项是
    \begin{tasks}(2)
        \task Grandient boosting
        \task cart树
        \task 随机森林
        \task AdaBoost
    \end{tasks}
\end{choice}


\begin{choice}{\;ACD\;}[]
    我国的  固定信令系统N0.7信令网采用的级别的分类包括
    \begin{tasks}(2)
        \task 低级信令转接点LTSP，汇接所属的SP
        \task 信令转接点STP
        \task 高级信令转接点HSTP，汇接LSTP和所属的SP
        \task 信令点SP
    \end{tasks}
\end{choice}


\begin{choice}{\;ABC\;}[]
    NR小区全局标识符（NCGI）由哪几部分组成
    \begin{tasks}(4)
        \task PLMN
        \task gNB ID
        \task Cell ID
        \task TAC
    \end{tasks}
\end{choice}

\begin{choice}{\;ABCD\;}[]
    在5G网络优化数据分析中，属于机器学习流程的选项是
    \begin{tasks}(2)
        \task 数据收集
        \task 数据建模
        \task 数据清洗
        \task 特征工程
    \end{tasks}
\end{choice}

\begin{choice}{\;ABCD\;}[]
    消费互联网和产业互联网的区别主要体现在哪些方面
    \begin{tasks}(2)
        \task 增长速度
        \task 市场结构
        \task 市场主角
        \task 服务对象
    \end{tasks}
\end{choice}

\begin{choice}{\;ABD\;}[]
    工业互联网平台作为工业全要素链接的枢纽与工业资源配置的核心，在工业互联网体系架构中具有至关重要的地位，是面向制造业什么的需求，构建基于海量数据采集、汇聚、分析的服务体系。
    \begin{tasks}(4)
        \task 智能化
        \task 数字化
        \task 信息化
        \task 网络化
    \end{tasks}
\end{choice}

\begin{choice}{\;ABCD\;}[]
    5G核心网的主要特征包括
    \begin{tasks}(4)
        \task 服务化架构
        \task 网络设备虚拟化
        \task 支持边缘计算
        \task 支持网络切片
    \end{tasks}
\end{choice}

\begin{choice}{\;CD\;}[]
    对于大唐5G设备EMB6216，5G基站开通网元布配的内容说法正确的是
    \begin{tasks}(1)
        \task 电源板卡在5槽位
        \task 风扇板卡无需配置
        \task  AAU光口的速率要与基带板光口的速率相同
        \task 基带板可以放在1槽位上
    \end{tasks}
\end{choice}

\begin{choice}{\;BC\;}[]
    6G网络架构中的四层包括
    \begin{tasks}(2)
        \task 智能功能层
        \task 资源与算力层
        \task 开放使能层
        \task 传输与接入层
    \end{tasks}
\end{choice}

\begin{choice}{\;ABCD\;}[]
    5G系统中网元AMF的功能包括
    \begin{tasks}(1)
        \task 承载管理功能，包括专用承载建立过程
        \task NAS移动性管理
        \task 用户鉴权及密钥管理
        \task NAS信令及信令的加密和完整性保护
    \end{tasks}
\end{choice}

\begin{choice}{\;ABC\;}[]
    SBA与传统核心网架构的区别在于
    \begin{tasks}(2)
        \task 移动性管理与会话管理解耦
        \task 控制面与媒体面分离
        \task 各种接入  方式都通过统一的机制接入网络
        \task 采用点对点通讯方式
    \end{tasks}
\end{choice}

\begin{choice}{\;BCD\;}[]
    管理会计理论中生产模式大致分为哪些类型
    \begin{tasks}(4)
        \task 迭代法
        \task 分批法
        \task 品种法
        \task 分步法
    \end{tasks}
\end{choice}

\begin{choice}{\;AD\;}[]
    对于eMBB场景三种NR典型帧结构，以下描述正确的是
    \begin{tasks}(2)
        \task 下行容量：2.5ms单 > 2ms单 > 2.5ms双
        \task 下行容量：2.5ms双 > 2ms单 > 2.5ms单
        \task 上行容量：2.5ms单 > 2ms单 > 2.5ms双
        \task 上行容量：2.5ms双 > 2ms单 > 2.5ms单
    \end{tasks}
\end{choice}

\begin{choice}{\;ABC\;}[]
    5G基站设备EMB6216开通时，传输不通需要排查的内容包括
    \begin{tasks}(1)
        \task SCTP链SCTP链路状态显示“驱动和高层建立成功”，说明IP路由层已通，注意检查基站和核心网的高层对接参数
        \task 检查基站侧传输光口下的指示灯状态，出现慢闪说明正常，出现闪烁和常亮交替，说明传输侧数据未配置完成
        \task SCTP链路状态显示“驱动和高层配置成功”，ARP状态为“学习中”说明数据链路层无问题，此时需要重点检查IP网络层参数
        \task SCTP链SCTP链路状态显示“未建”，说明数据链路层故障，需要排查VLAN设备是否正确
    \end{tasks}
\end{choice}

\begin{choice}{\;ABD\;}[]
    MIB消息携带了5G系统的有关信息，以下内容属于MIB携带内容的是
    \begin{tasks}(4)
        \task 子载波间隔
        \task 小区禁止信息
        \task PLMN
        \task 系统帧号
    \end{tasks}
\end{choice}

\begin{choice}{\;BD\;}[]
    作业成本法的优点包括
    \begin{tasks}(2)
        \task 利于管理控制
        \task 有助于改进成本控制
        \task 开发和维护费用低
        \task 可以获得更准确的产品和产品线成本
    \end{tasks}
\end{choice}


\begin{choice}{\;BD\;}[]
    目前C-V2X正在向交通安全和效率类应用发展，并逐步向支持实现自动驾驶的协同服务类应用演进，以下属于交通效率类应用的包括
    \begin{tasks}(2)
        \task 车辆编队
        \task 车速引导
        \task 远程遥控驾驶
        \task 紧急呼叫
    \end{tasks}
\end{choice}


\begin{choice}{\;BCD\;}[]
    技术创新方式包括哪几种
    \begin{tasks}(4)
        \task 自主创新
        \task 原始创新
        \task 集成创新
        \task 改进创新
    \end{tasks}
\end{choice}


\begin{choice}{\;AB\;}[]
    5G移动通信系统中，URLLC场景下的业务必须保证
    \begin{tasks}(4)
        \task 低时延
        \task 高可靠
        \task 高速率
        \task 低功耗
    \end{tasks}
\end{choice}


\begin{choice}{\;ABC\;}[]
    5G NR中，毫米波支持的信道带宽是
    \begin{tasks}(4)
        \task 50MHZ
        \task 100MHZ
        \task 200Mhz
        \task 10MHZ
    \end{tasks}
\end{choice}


\begin{choice}{\;BCD\;}[]
    移动通信系统中，关于信号功率的计算，以下计算公式正确的是
    \begin{tasks}(2)
        \task $XdBm+Ydbm=(X+Y)dBm$
        \task $XdB-YdB=(X-Y)dB$
        \task $XdBm+YdB=(X+Y)dBm$
        \task $XdB+YdB=(X+Y)dB$
    \end{tasks}
\end{choice}


\begin{choice}{\;ABC\;}[]
    对于R15版本的5G NR，以下说法正确的是
    \begin{tasks}(1)
        \task 当$\mu =1$时，SCS=30khz，可配置普通CP，此时一个无线帧中包含280个符号
        \task 当$\mu =2$时，SCS=60khz，可配置扩展CP，此时一个无线帧中包含480个符号
        \task 当$\mu =2$时，SCS=60khz，可配置普通CP，此时一个无线帧中包含560个符号
        \task 当$\mu =1$时，SCS=30khz，可配置扩展CP，此时一个无线帧中包含240个符号
    \end{tasks}
\end{choice}


\begin{choice}{\;AD\;}[]
    关于5G架构选项2，以下描述错误的是
    \begin{tasks}(2)
        \task 它的核心网是EPC
        \task 它的无线接入技术使用NR
        \task 它是独立架构
        \task 它通过eUTRAN实现用户控制面连接
    \end{tasks}
\end{choice}


\section{判断题（共30小题，共60分）}

\begin{choice}{\;正确\;}[]
    边缘计算：适用于实时，短周期数据，本地决策场景。
\end{choice}


% \begin{choice}{\;错误\;}[]
%     RRC连接建立失败会触发RRC重建
% \end{choice}

\begin{choice}{\;正确\;}[]
    移动通信系统中，多普勒频移的大小和无线电波的波长有很大的关系。
\end{choice}

\begin{choice}{\;正确\;}[]
    分离承载是指可以由MeNB或SeNB与GW对接，再通过Xn/Xx接口数据分离。
\end{choice}

\begin{choice}{\;正确\;}[]
    模拟信号和数字信号的区别可根据波形是否连续来区别确定。
\end{choice}

\begin{choice}{\;正确\;}[]
    静止卫星的轨道必然处于赤道平面上。

\end{choice}


\begin{choice}{\;错误\;}[]
    CAN总线远程帧不存在数据场，所以DLC的数值应永远被设置为0。

\end{choice}

\begin{choice}{\;错误\;}[]
    5G基本时间单位$T_c$和LTE基本时间单位$T_s$是一样的。

\end{choice}

\begin{choice}{\;正确\;}[]

    在接收端，数据要经过解调，解扩，解码等过程进行传输。
\end{choice}


\begin{choice}{\;正确\;}[]
    以技术为主，以研发部为中心的传统产品开发模式，是面向技术先进性而进行的产品开发，不是以客户为中心的产品开发模式。

\end{choice}

\begin{choice}{\;错误\;}[]
    对于过覆盖优化，我们通常采取上抬下倾角式，此举一来可以抑制过覆盖，同样可以缓解灯下黑的状况。

\end{choice}


\begin{choice}{\;错误\;}[]
    基尼指数越大，样本的确定性越大。

\end{choice}


\begin{choice}{\;正确\;}[]
    SLAM是Simultaneous\;localization\;and\;mapping缩写，意为同步定位与建图，主要用于解决机器人在未知环境运动时的定位与地图构建问题。

\end{choice}


\begin{choice}{\;正确\;}[]
    GPS避雷器安装在GPS射频馈线进入馈线窗后1m处。

\end{choice}


\begin{choice}{\;正确\;}[]
    人工智能在5G+无线接入网方面可以提供网络的管理和优化。

\end{choice}


\begin{choice}{\;正确\;}[]
    工业互联网是跨界融合的系统性工程，是与工业生产紧密相关的新型网络基础设施。

\end{choice}

% \begin{choice}{\;正确\;}[]
%     以技术为主，以研发部为中心的传统产品开发模式，是面向技术先进性而进行的产品开发，不是以客户为中心的产品开发模式

% \end{choice}


\begin{choice}{\;错误\;}[]
    C-V2X相对于ADAS更成熟更早商用。

\end{choice}


\begin{choice}{\;正确\;}[]
    移动通信系统中，扩频码会改变信号的带宽，扰码不会改变信号的带宽。

\end{choice}


\begin{choice}{\;错误\;}[]
    控制与转发分离：控制面进行统一的策略控制，转发面更关注转发功能的实现，以软件来驱动，实现了软硬件的解耦。

\end{choice}


\begin{choice}{\;正确\;}[]
    一般产品都具备技术属性和经济属性，不以人的意志为转移。

\end{choice}


\begin{choice}{\;正确\;}[]
    MCG\;split\;bearer的特性：在双连接里，承载无线协议在MgNB被分裂，并且承载既属于MCG，也属于SCG。

\end{choice}


\begin{choice}{\;正确\;}[]
    C-V2X的接口只有PC5和Uu口。

\end{choice}

\begin{choice}{\;错误\;}[]
    5G的RRC消息在传送前还需要经过UDP/IP的封装过程。

\end{choice}


\begin{choice}{\;正确\;}[]
    缩小宏站的覆盖距离，不一定能提升覆盖性能。

\end{choice}

\begin{choice}{\;错误\;}[]
    5G信令中，UE建立RRC连接建立后，RRC的状态可能是RRC CONNECTED或RRC INACTIVE状态。

\end{choice}

\begin{choice}{\;错误\;}[]
    5G的RRC消息在传送前还需要经过UDP/IP封装过程，才能过通过DRB经由LTE传送给UE。

\end{choice}

\begin{choice}{\;正确\;}[]
    电信网按照业务性质分类，可以分为业务网、传输网、支撑网。

\end{choice}

\begin{choice}{\;正确\;}[]
    5G多网络融合技术可以考虑分布式和集中式两种结构。

\end{choice}

\begin{choice}{\;错误\;}[]
    必须使用静止卫星才能实现卫星通信。

\end{choice}

\begin{choice}{\;正确\;}[]
    在5G NSA组网下，UE和网络可能建立MCG bearer，SCG bearer。

\end{choice}

\begin{choice}{\;错误\;}[]
    产品的成本及售价，可以在设计之后再行设定。

\end{choice}





